\section*{Chapter \ref{ch:g6:if-else} --- Decisions: \texttt{if} and \texttt{else}}

\begin{solution}{exo:g6:if-else:1}
\code{True} and \code{True}: the whole number 4 and the decimal number 4.0 are
equal, even though their types differ.
\end{solution}

\begin{solution}{exo:g6:if-else:2}
\code{medium}: 7 is not more than 10, so the first \omterm{def:g4:algorithms-repeat-and-decide:condition}{condition} is false; it is
more than 5, so the \code{elif} \omterm{def:g6:if-else:if}{branch} runs, and the \code{else} is skipped.
\end{solution}

\begin{solution}{exo:g6:if-else:3}
\code{a and b} is \code{False}, so \code{not (a and b)} is \code{True}.
\code{not a} is \code{False} and \code{b} is \code{False}, so \code{not a or b}
is \code{False}: \code{not} applies only to \code{a}.
\end{solution}

\begin{solution}{exo:g6:if-else:4}
\omcode{code/grade-6/04-if-else/exercises/sign.py}{1}{5}{Positive or
negative.}\label{lst:g6:if-else:sol-sign}

Zero goes to the first \omterm{def:g6:if-else:if}{branch} because the \omterm{def:g6:if-else:comparison}{comparison} is \code{>=}.
\end{solution}

\begin{solution}{exo:g6:if-else:5}
\code{x == 3} is a \omterm{def:g6:if-else:comparison}{comparison}, a question: it computes \code{False}, since \code{x}
is 1, and the answer is thrown away. Nothing changes \code{x}, and nothing is
wrong in the language, so there is no \omterm{def:g6:first-python-program:error}{error message}. To give \code{x} the \omterm{def:g6:variables-and-input:value}{value} 3,
the line must be an \omterm{def:g6:variables-and-input:assignment}{assignment}, with one \code{=}: \code{x = 3}. It is the mirror of
the slip of \cref{out:g6:if-else:assign}, where one \code{=} stood where
\code{==} was needed.
\end{solution}

\begin{solution}{exo:g6:if-else:6}
Anyone older than 65 is also older than 12, so the first \omterm{def:g4:algorithms-repeat-and-decide:condition}{condition} is already
true and the \code{elif} is never checked. The narrower case must come first:

\omcode{code/grade-6/04-if-else/exercises/senior.py}{1}{7}{The senior case
first.}\label{lst:g6:if-else:sol-senior}
\end{solution}

\begin{solution}{exo:g6:if-else:7}
The \omterm{def:g4:programs-and-operating-system:program}{program} keeps the largest number seen so far in a \omterm{def:g6:variables-and-input:value}{variable} and compares
each new number with it.

\omcode{code/grade-6/04-if-else/exercises/largest.py}{1}{9}{The largest of
three.}\label{lst:g6:if-else:sol-largest}

For 4, 9 and 2 it prints 9.0. It should be tested with the largest number in
each of the three places, and with two equal numbers.
\end{solution}

\begin{solution}{exo:g6:if-else:8}
The highest band is tested first, so each \code{elif} only sees the marks below
the bands above it.

\omcode{code/grade-6/04-if-else/exercises/grade.py}{1}{10}{A grade
letter.}\label{lst:g6:if-else:sol-grade}
\end{solution}

\begin{solution}{exo:g6:if-else:9}
One \omterm{def:g6:variables-and-input:value}{value} inside each band and the \omterm{def:g6:variables-and-input:value}{values} on both sides of each boundary: $-5$
(ice), $-1$ and $0$ (ice, then cool), $19$ and $20$ (cool, then warm), $30$ and
$31$ (warm, then hot). The boundary 30 belongs to \code{warm} because the test is
\code{<=}.
\end{solution}

\begin{solution}{exo:g6:if-else:10}
\omcode{code/grade-6/04-if-else/exercises/password.py}{1}{5}{A password-length
check.}\label{lst:g6:if-else:sol-password}

For \code{abc} it says that 9 more characters are needed. It should be tested
with 11, 12 and 13 characters.
\end{solution}

\begin{solution}{exo:g6:if-else:11}
They always agree. If \code{a} is false, the first way goes to the outer
\code{else} (no), and \code{a and b} is false (no). If \code{a} is true, the
first way answers as \code{b} decides, and \code{a and b} is true exactly when
\code{b} is. The \omterm{def:g4:programs-and-operating-system:program}{program} prints \code{yes yes} only for \code{True True}, and
\code{no no} in the three other cases.
\end{solution}

\begin{solution}{exo:g6:if-else:12}
The \omterm{def:g4:programs-and-operating-system:program}{program} tests only ``divisible by 4'' and forgets the exception for the
centuries. The fix is the \omterm{def:g4:algorithms-repeat-and-decide:condition}{condition} of \cref{lst:g6:if-else:leap}:
\texttt{year \% 4 == 0 and (year \% 100 != 0 or year \% 400 == 0)}. The year 1900
shows the fix: the old \omterm{def:g4:programs-and-operating-system:program}{program} says leap year, the new one says not. The year
2000 must also be tried, to check that the exception of the exception still
works: it is a leap year.
\end{solution}

\begin{solution}{pb:g6:if-else:1}
\textbf{1.} 2024 is divisible by 4 and does not end a century: a leap year.
2023 is not divisible by 4: not a leap year.

\textbf{2.} 2100 ends a century and is not divisible by 400: not a leap year.
2000 is divisible by 400: a leap year. 1900, like 2100: not a leap year.

\textbf{3.} 366 days. Four years with one leap year: $3 \times 365 + 366 =
1461$ days.

\textbf{4.} Both need \texttt{year \% 4 == 0}. For 2024 the bracket is true
because \texttt{year \% 100 != 0} (2024 does not end a century); for 2000 it is
true because \texttt{year \% 400 == 0}.

\textbf{5.} The \omterm{def:g4:programs-and-operating-system:program}{program} agrees: leap years for 2024 and 2000, not for 2100
and 1900.

\textbf{6.} and \textbf{7.} The \omterm{def:g4:programs-and-operating-system:program}{program} below does both counts:

\omcode{code/grade-6/04-if-else/listings/count_leap.py}{1}{14}{Counting leap
years.}\label{lst:g6:if-else:count-leap}

\omoutput{code/grade-6/04-if-else/outputs/count_leap.txt}{49 and
97.}\label{out:g6:if-else:count-leap}

From 1901 to 2100 there are 49 leap years: the 50 multiples of 4 from 1904 to
2100, minus 2100. From 2001 to 2400 there are 97, in twice as many years.

\textbf{8.} The rule depends only on the remainder of the year divided by 400:
the year $y + 400$ has the same remainders by 4, 100 and 400 as $y$. So any 400
years in a row have the same leap years, moved by a multiple of 400.

\textbf{9.} In 400 years, 100 years are divisible by 4. Four of them end a
century, and one of those four is divisible by 400. So $100 - 4 + 1 = 97$ leap
years.

\textbf{10.} $400 \times 365 + 97 = 146{,}097$ days.

\textbf{11.} $146{,}097 \div 400 = 365.2425$ days. Every fourth year a leap year
would give $365.25$: the century rule takes away 3 days in 400 years, and makes
the average year $0.0075$ days shorter, much closer to the time the Earth takes
to go round the Sun.
% ledger: g6:gregorian-mean

\textbf{12.} \textbf{Result.} There are 97 leap years in every 400 years, so
400 years have 146,097 days, and the average year of the calendar lasts 365.2425
days.
\end{solution}
